\chapter{Conclusion}
Dynamic positioning behaviour of ships in offshore has traditionally been predicted using quasi-static force balances.
The industry is now transitioning to time-domain simulations, which better capture transients and footprint behavior.
These simulations are computationally costly and still subject to model error, and perfect knowledge of the simulated
system is rarely available. Quantifying aleatoric and epistemic uncertainty in a DP system remains an open research
question, and therefore the central topic of this thesis.

Based on the gap in the research on dynamic positioning simulation, a research question was formulated as follows:

\emph{How can a ship's short-term dynamic positioning behaviour efficiently be forecast while considering environmental and system uncertainty?}

To answer this question, a set of research goals were set. The first two of these are covered by the literature review
and scoping study. 

Current deterministic simulation frameworks are unsuitable for uncertainty quantification, as the many simulations
required for Monte-Carlo UQ would be intractable for operational decision making. For this reason, surrogate approaches
were investigated. Across the literature, uncertainty is commonly separated into aleatoric and epistemic components.
Most surrogate pipelines capture stochastic dynamics in the surrogate model, after which a UQ method is applied to
characterize model uncertainty.

Several UQ methods were reviewed, with polynomial chaos expansions (PCE) and machine learning approaches as the most
prominent. PCE is widely used for UQ but becomes inefficient for long time-domain simulations and suffers from the curse
of dimensionality. Machine learning methods are more recent but offer greater flexibility for probabilistic prediction.
Among these, Bayesian neural networks and deep ensembles are well aligned with the need to quantify uncertainty in
surrogate DP models.

For modelling the dynamics, two main approaches stand out. Operator learning, through DeepONet or Koopman methods, offers
computationally inexpensive inference and enables fast propagation of stochastic signals through the system dynamics.
These approaches can be paired with epistemic UQ methods. Between the two, DeepONet appears more suitable for the task
at hand. Neural stochastic differential equations provide an alternative by directly modelling stochastic dynamics with
learned drift and diffusion terms. In the pilot tests on a toy DP model, neural SDEs performed better than DeepONet.

Through the literature review and initial tests, the first two research goals have been completed.

\section{Thesis Planning}
The literature review and initial tests have provided a solid backbone for the next stages of the thesis.

First, frameworks for working with the neural SDEs will be investigated. In particular, the use of Diffrax will be
tested. Using the most suitable approach from this study, the toy model problem will be fully trained to confirm the
ability of neural SDEs to model DP behaviour. 

LatentSDE vs SDE-GAN?
Deep ensembles for epistemic uncertainty quantification
More complex DP model - Convert to MSS example to python, generate a whole lot of training data
Generalisation of NeuralSDE with different model parameters as inputs
Transient thruster failures